\documentclass[notitlepage,onecolumn,showpacs,11pt]{revtex4-1}
\usepackage{amssymb}
\usepackage{graphicx}
\usepackage{amsmath}
\usepackage{float}
\usepackage{epstopdf}
\usepackage{times}
\usepackage{color}
\usepackage{subfigure}
\usepackage{setspace}
\usepackage{bm}% bold math

\newcommand{\bsym}[1]{\boldsymbol{#1}}
\newcommand{\mbf}[1]{\mathbf{#1}}
\newcommand{\grad}{\boldsymbol{\nabla}}
\newcommand{\Tr}{\text{Tr}}

\begin{document}

\title{Bond-disordered non-collinear magnets: Perturbative renormalization}

\date{\today}

\maketitle

We consider the appropriate continuum modelling for bond disordered
non-collinear antiferromagnets and investigate the role of disorder in
the fate of long range order in these systems.


\section{NL$\sigma$-model for non-collinear antiferromagnets}


In non-collinear magnets the original symmetry of the spins is fully broken, 
e.g. in the triangular lattice heisenberg antiferromagnet the broken symmetry 
is 
\begin{equation}
    O(3)/Z_2\simeq SO(3)
\end{equation}
.As shown by Azaria et. al \cite{azaria-93}, for a complete description of the 
low energy excitation one has to also consider the planar nature of the ground 
state and distinguish between the three Goldstone modes of the problem. In the 
following we follow a microscopic analysis \cite{dombre-88} to obtain the 
coarse-grained effective hamiltonian of the triangular lattice antiferromagnet 
with continuous symmetry. Later, we shall extend this scheme to consdier the 
effect of bond disorder.. 

We consider a non-collinear ordered state. The Goldstone modes shall be excited 
by the generators of the broken symmetry group, $G/H$. With the trinagular
lattice in mind, our continuum modelling involves considering plaquettes of 
elementary upward triangles, finding the low energy contributions for 
departure from the $120^\circ$ order and ultimately coarse-graining within 
each elementary triangle. We start by denoting the three sublattice 
orientaions of the $120^\circ$ order, $\mbf{n}_A,\mbf{n}_B$ and 
$\mbf{n}_C = -(\mbf{n}_A+\mbf{n}_B)$ with spin-excitations generated by 
$\mbf{S}_\alpha=U\mbf{n}_\alpha$, where $U\in G/H$. We consider the
hamiltonian within an elementary triangular plaquette.
\begin{equation}
    \begin{aligned}
        H = J \left(\mbf{S}_A(\mbf{x}_A)\cdot\mbf{S}_B(\mbf{x}_B)
        +\mbf{S}_B(\mbf{x}_B)\cdot\mbf{S}_C(\mbf{x}_C)
        +\mbf{S}_C(\mbf{x}_C)\cdot\mbf{S}_A(\mbf{x}_A)
        \right) 
    \end{aligned}
\end{equation}
By performing a gradient expansion coupled with a coarse-graining of the 
spatial co-ordinate to capture the low energy pieces, we arrive at the 
following expression.
\begin{equation}
    \begin{aligned}
        H &= JS^2a^2 \left(\mbf{n}_A U^{-1}\cdot\grad^2U\mbf{n}_B
        -\mbf{n}_B U^{-1}\cdot\grad^2U(\mbf{n}_A+\mbf{n}_B)
        -\mbf{n}_A U^{-1}\cdot\grad^2U(\mbf{n}_A+\mbf{n}_B)
        \right)\\
        &= JS^2a^2 \left(\mbf{n}_A U^{-1}\cdot\grad^2U\mbf{n}_B
        -(\mbf{n}_A+\mbf{n}_B)U^{-1}\cdot\grad^2U(\mbf{n}_A+\mbf{n}_B)
        \right) 
    \end{aligned}
\end{equation}
The single gradient term cancels from the expansion. Now, a transformation 
to the orthonormal basis, $\mbf{e}_1=\left(\mbf{n}_A+\mbf{n}_B\right)$ and 
$\mbf{e}_2=\frac{1}{\sqrt{3}}\left(\mbf{n}_A-\mbf{n}_B\right)$ re-arranges 
the terms of the hamiltonian in a convenient form.
\begin{equation}
    \begin{aligned}
        H &= -\frac{3JS^2a^2}{2} \left(\mbf{e}_1 U^{-1}\cdot\grad^2U\mbf{e}_1
        +\mbf{e}_2 U^{-1}\cdot\grad^2U\mbf{e}_2
        \right)
    \end{aligned}
\end{equation}
An extra $O(2)$ symmetry within the $(\mbf{e}_1,\mbf{e}_2)$ vector space is 
immediately apparent. The lowest non-trivial representation, $U$ for $SO(3)$
is three-dimensional. Therefore, by introducing,

\begin{itemize}
    \item A third orthonormal vector $\mbf{e}_3=\mbf{e}_1\times\mbf{e}_{2}$.
    \item A projector $P\sim\text{diag}\left(1,1,0\right)$ to the 
        $(\mbf{e}_1,\mbf{e}_2)$ plane.
    \item Defining $S_0=(\mbf{e}_1,\mbf{e}_2,\mbf{e}_3)$ and performing a 
        gauge transformation $U = US_0$ we can re-write,
\end{itemize}

\begin{equation}
    \begin{aligned}
        H &= -\frac{3JS^2a^2}{2} 
        \Tr\left(S_0^{-1}P U^{-1}\grad^2U PS_0\right)\\
        &= -\frac{3JS^2a^2}{2} \Tr\left(PU^{-1}\grad^2U\right)
    \end{aligned}
\end{equation}
However, for a proper RG treatment of the action it is necessary to consider 
a separate copuling constant \cite{azaria-93}. 
This introduces a second coupling to the problem with the complete hamiltonian 
now given by,
\begin{equation}
    \begin{aligned}
        H &= -p_1\left(\mbf{e}_1 U^{-1}\cdot\grad^2U\mbf{e}_1
        +\mbf{e}_2 U^{-1}\cdot\grad^2U\mbf{e}_2
        \right)
        -p_2\left(\mbf{e}_3 U^{-1}\cdot\grad^2U\mbf{e}_3\right)\\
        &=-\Tr\left(P(U^{-1}\grad U)^2\right)
    \end{aligned}
\end{equation}
with gauge transformation $U = U (\mbf{e}_1,\mbf{e}_2,\mbf{e}_3)$ generating 
the full $SO(3)$ Goldstone modes and the diagonal coupling matrix 
$P=\text{diag}(p_1,p_1,p_2)$ ensuring the $O(3)\times O(2)/O(2)$ nature of 
the exciatations. Given $\pi_1(SO(3))=Z_2$, the order parameter $U$ can have 
$Z_2$ point defects which are gapped out at zero temperature. Therefor the 
quantum action only has terms with second time derivatives and the full 
Goldstone mode action is therefore given by \cite{dombre-88,azaria-92}.
\begin{equation}
    \begin{aligned}
        S &=-\int d\tau\int d^dx 
        \left[
        \Tr\left(P_\tau(U^{-1}\partial_\tau U)^2\right)
        +
        \Tr\left(P_x(U^{-1}\grad U)^2\right)
        \right]
    \end{aligned}
\end{equation}
where the two distinct diagonal matrices, $P_\tau$ and $P_x$ presupposes the 
existence of three different spin-wave velocities of the problem as also 
obtained from a spin-wave expansion around the ordered state.


\section{Bond-disordered non-collinear antiferromagnet}

First, we consider an isolated impurity within one of the elementary plaquettes
and we find that the additional contribution from a bond impurity across 
sublattices $AB$, $J_{AB} = J + \delta J$ is given by,
\begin{equation}
    \begin{aligned}
        \delta H &= \delta J \ \mbf{S}_A(\mbf{x}_A)\cdot\mbf{S}_B(\mbf{x}_B)
    \end{aligned}
\end{equation}
. Unlike the bulk, the energy cost from Goldstone mode excitations now has
contribution from the uncompensated single gradient term. Like before, we 
use the convenient orthonormal basis transformation and find the impurity 
contribution to be of the following form.
\begin{equation}
    \begin{aligned}
        \delta H &= \delta JS^2a \ \mbf{n}_AU^{-1}\cdot(\mbf{u}_{AB}\cdot\grad)U\mbf{n}_B\\
        &= -\frac{\sqrt{3}\delta JS^2a}{4}\left(
        \mbf{e}_2 U^{-1}\cdot(\mbf{u}_{AB}\cdot\grad)U \mbf{e}_1
        -\mbf{e}_1 U^{-1}\cdot(\mbf{u}_{AB}\cdot\grad)U\mbf{e}_2
        \right)
    \end{aligned}
\end{equation}
where we have exploited the orthonormality of the basis to drop the terms 
$\mbf{e}_1U^{-1}\grad U\mbf{e}_1$ and $\mbf{e}_2 U^{-1}\grad U\mbf{e}_2$ which 
are zero. For a bond impurity across a bond orientation $\mbf{u}_\alpha$ 
the pertubation assumes the form,
\begin{equation}
    \begin{aligned}
        \delta H &= h^\alpha\mbf{u}_\alpha\cdot\Tr\left(
        P_xt^{3}
        \begin{pmatrix}
            \mbf{e}_1^T\\
            \mbf{e}_2^T\\
            \mbf{e}_3^T
        \end{pmatrix}
        U^{-1}
        \grad
        U
        (\mbf{e}_1,\mbf{e}_2,\mbf{e}_3)\right) \\
        &= h^\alpha\mbf{u}_\alpha
        \cdot\Tr\left(P_xt^3 U^{-1}\grad U\right)
    \end{aligned}
\end{equation}
. The same gauge choice and projector from the bulk action was used and the 
$SO(3)$ generator corresponding to the in plane rotation
\begin{equation}
    t^3=\begin{pmatrix}0&1&0\\-1&0&0\\0&0&0\end{pmatrix}
\end{equation}
makes an automatic appearance. This tells us that in the triangular lattice 
the local dipolar field arising due to bond disorder couples directly with 
the Goldstone mode associated with the in plane rotation. Now, it is 
straightforward to generalize to the case of spatially distributed bond 
impurities i.e. disordered bonds.
\begin{equation}
    \begin{aligned}
        H = \int d^dx\left[
            \Tr\left(P_x(U^{-1}\grad U)^2\right)
            +h(\mbf{x})^\alpha
            \mbf{u}_\alpha
            \cdot\Tr\left(P_xt^3(U^{-1}\grad U)\right)
        \right]
    \end{aligned}
\end{equation}

\section{Replica renormalization group}

We consider a gaussian distribution for the disordered couplings,
\begin{equation}
    \begin{aligned}
        \mathcal{P}\left[h\right]
        &=\int\mathcal{D}h\text{e}^
        {-\frac{1}{2\Delta^2}\int d^dx
        h^\alpha(\mbf{x})
        h^\alpha(\mbf{x})
        }
    \end{aligned}
\end{equation}
which allows us to use the replica trick and write down a homogeneous
action suitable for the bond-disordered non-collinear order, albeit, at
the expense of added replica indices. 
\begin{equation}
    \begin{aligned}
        S =& -\sum_r\int d\tau\int d^dx
        \Tr\left[P_\mu
            \left(U_r^{-1}\partial_\mu U_r\right)_{G}^2
            \right] \\
        &-\Delta^2
        \sum_{r,s}\int d\tau\int d\tau'\int d^dx
        \Tr\left[P_xt^3
            \left(U_r^{-1}\grad U_r\right)_G
            \right]
            \cdot
            \Tr\left[P_xt^3
            \left(U_s^{-1}\grad U_s\right)_G
            \right]
    \end{aligned}
\end{equation}
Where the concrete example for the triangular lattice Heisenberg 
antiferromagnet shall have 
$\mathcal{L}\left[G\right]=\mathcal{L}\left[SO(3)\right]$ and 
$t^3 \in \mathcal{L}\left[SO(2)\right]$. We perform a trivial rescaling 
to restore the space-time Lorentz symmetry and re-write conveniently,
\begin{equation}
    \begin{aligned}
        S =& -\sum_r\int d\tau\int d^dx
        \left(\frac{g_1}{2}\Tr\left[
            \left(U_r^{-1}\partial_\mu U_r\right)_{G/X}^2\right]
        +\frac{g_2}{2}\Tr\left[
            \left(U_r^{-1}\partial_\mu U_r\right)_{X}^2
            \right]
        \right)\\
        &-\frac{\Delta^2g_2^2}{4}
        \sum_{r,s}\int d\tau\int d\tau'\int d^dx
        \Tr\left[t^3
            \left(U_r^{-1}\grad U_r\right)_G
            \right]
            \cdot
            \Tr\left[t^3
            \left(U_s^{-1}\grad U_s\right)_G
            \right]
    \end{aligned}
\end{equation}
with $X=SO(2)$. The dimensional natures of the bulk and disordered couplings 
are immediately apparent,
\begin{equation}
    \begin{aligned}
        \left[g_{1,2}\right] &= d - 1\\
        \left[\Delta\right] &= (2-d)/2
    \end{aligned}
\end{equation}
The bulk coupling $g_{1,2}$ is relevant and consistent with the fact that the
homogeneous system is ordered. In this case $g_{1,2}$ has a trivial 
renormalization,
\begin{equation}
    \begin{aligned}
        Z_{g_{1,2}}g_{1,2}\Lambda^{d-1} &= g_B \\
        \implies \beta(g_{1,2})
        &=(1-d)\Lambda^{d-1}g_{1,2} + \mathcal{O}(1)
    \end{aligned}
\end{equation}
In $d=1$, the bulk coupling becomes marginal and disorder is relevant. 
Disorder shifts the wilsonian fixed point of the bulk action. However, in
the following we shall focus on $d>=2$ and try to track the renormalization
of the disorder coupling with the help of a perturbative expansion of the
Goldstone mode action. 

We use the exponential representation for the full symmetry group, 
$U=e^{\bsym{\xi}\cdot\mbf{t}}$ with $t^a\in \mathcal{L}\left[G\right]$
and $\xi_a$ are the associated Goldstone modes. We identified the
in-plane Goldstone mode $\in X$ to be special and there is an easy 
separation between the respective generators of $\mathcal{L}(G/X)$ and 
$\mathcal{L}(X)$. Using the notation,
\begin{equation}
    \begin{aligned}
        t^a &\in \mathcal{L}(G/X) \\
        t^\alpha &\in \mathcal{L}(X)
    \end{aligned}
\end{equation}
we have the commutation relations,
\begin{equation}
    \begin{aligned}
        \left[t^\alpha,t^\beta\right] 
        &= f_{\alpha\beta\gamma}t^\gamma \\
        \left[t^\alpha,t^b\right] 
        &= f_{\alpha bc}t^c \\
        \left[t^a,t^b\right] 
        &= f_{ab\gamma}t^\gamma  \text{ \ and \ }\\
        \Tr\left[t^at^b\right] &= -\delta^{ab}
    \end{aligned}
\end{equation}
With the Lie algebra defined it is now straightforward to expand the 
terms of the action in terms of the Goldstone modes.
\begin{equation}
    \begin{aligned}
        -\Tr\left[\left(U^{-1}\partial_\mu U\right)^2_{G/X}\right]
        &=(\partial_\mu\xi_a)^2 
        + \frac{1}{3}\left(f_{ab\gamma}\xi_a\partial_\mu\xi_b\right)^2
        + \hdots\\
        -\Tr\left[\left(U^{-1}\partial_\mu U\right)^2_{X}\right]
        &=(\partial_\mu\xi_\alpha)^2\\
        -\Tr\left[
            t^\alpha\left(U^{-1}\partial_\mu U\right)_G
            \right]
        &= \partial_\mu\xi_\alpha
            +\frac{K_{bc\alpha}}{2!}\left(\partial_\mu\xi_b\right)\xi_c
            +\frac{K_{bcd\alpha}}{3!}\left(\partial_\mu\xi_b\right)
            \xi_c\xi_d
            +\hdots
    \end{aligned}
\end{equation}
where we have used the fact that $X=SO(2)$ is an abelian group and used
a compact notation $K_{bc\alpha} = \left[t^b,t^c\right]_\alpha=f_{bc\alpha}$, 
$K_{bcd\alpha} = \left[\left[t^b,t^c\right],t^d\right]_\alpha
=f_{bcx}f_{xd\alpha}$ and so on. The explicit form of the action therefore is.
\begin{equation}
    \begin{aligned}
        S =& \sum_r\int d\tau\int d^dx
        \left(\frac{g_1}{2}\left[
           (\partial_\mu\xi^r_a)^2 
        + \frac{1}{3}\left(f_{ab\gamma}\xi^r_a\partial_\mu\xi^r_b\right)^2
        + \hdots\right]
        +\frac{g_2}{2}\left[
            (\partial_\mu\xi^r_\alpha)^2            
            \right]
        \right)\\
        &-\frac{\Delta^2g_2^2}{4}
        \int d^dx
        \left[\sum_r\int d\tau
            \left(
            \partial_\mu\xi^r_\alpha
            +\frac{K_{bc\alpha}}{2!}\left(\partial_\mu\xi^r_b\right)\xi^r_c
            +\frac{K_{bcd\alpha}}{3!}\left(\partial_\mu\xi^r_b\right)
            \xi^r_c\xi^r_d
            +\hdots
            \right)
            \right]^2
    \end{aligned}
\end{equation}
It is interesting to note the nature of infrared divergence of the theory, 
for the $\xi^r_\alpha \in SO(2)$ bosons, driven by the quenched disorder 
term of the action. One way to regulate this divergence is to introduce an 
external field term which should ultimately be taken to zero
\cite{brezin-76}. Introducing field strength, coupling constant and 
disorder strength renormalization constants, $Z$, $Z_{g_{1,2}}$ and 
$Z_\Delta$ we write down the full renormalizable replica theory for 
our problem.
\begin{equation}
    \begin{aligned}
        S_{ren} =& \sum_r\int d\tau\int d^dx
        \left(Z_{g_1}\Lambda^{d-1}\frac{g_1}{2}
        \left[
            Z(\partial_\mu\xi^r_a)^2 
        + \frac{1}{3}\left(Zf_{ab\gamma}\xi^r_a\partial_\mu\xi^r_b\right)^2
        + \hdots\right]
        +Z_{g_2}\Lambda^{d-1}\frac{g_2}{2}
        \left[Z
            (\partial_\mu\xi^r_\alpha)^2            
            \right]
        \right)\\
        &-Z_{\Delta}^2Z^2_{g_2}\Lambda^d\frac{\Delta^2g_2^2}{4}
        \int d^dx
        \left[\sum_r\int d\tau
            \left(
            \sqrt{Z}\partial_\mu\xi^r_\alpha
            +Z\frac{K_{bc\alpha}}{2!}\left(\partial_\mu\xi^r_b\right)\xi^r_c
            +Z^{3/2}\frac{K_{bcd\alpha}}{3!}\left(\partial_\mu\xi^r_b\right)
            \xi^r_c\xi^r_d
            +\hdots
            \right)
            \right]^2\\
            &-Z^{-1/2}\Lambda^{d-1}\frac{g_2}{2}
            m^2\sum_r\int d\tau\int d^dx
            \left[
                1-\frac{Z}{2!}\left(\xi^r_\alpha\right)^2
                +\frac{Z^2}{4!}\Tr\left[t^\alpha t^\alpha
                t^\alpha t^\alpha\right]
                \left(\xi^r_\alpha\right)^4
                \hdots
                \right]
    \end{aligned}
\end{equation}
We consider the Greens function for the $\xi^r_\alpha$ Goldstone modes
which gives rise to the logarithmic divergence of the theory.
\begin{equation}
    \begin{aligned}
        G^{rs}_\alpha(\omega,\mbf{k})
        &=
        \frac{1}{Z_{g_2}Z\Lambda^{d-1}}
        \frac{1}{g_2}
        \frac{\delta^{rs}}{\omega^2+\mbf{k}^2+m^2_0}
        +
        \frac{Z_\Delta^2\Lambda^{2-d}}{Z}
        \frac{\Delta^2}{2}
        \frac{\mbf{k}^2}
        {\left(\omega^2+\mbf{k}^2+m_0^2\right)^2}
        2\pi\delta(\omega)
    \end{aligned}
\end{equation}
with,
\begin{equation}
    \begin{aligned}
        m^2_0 = Z^{1/2}\Lambda^{d-1}\frac{g_2}{2}m^2
    \end{aligned}
\end{equation}
The two-point function for $\xi^r_\alpha$ bosons up to one loop is,
\begin{equation}
    \begin{aligned}
        \Gamma^{(2)}_{rs}\left(\omega,\mbf{k}\right)
        =&Z_{g_2}Z\Lambda^{d-1}g_2(\omega^2+\mbf{k}^2)\delta_{rs}
        -
        Z_{\Delta}^2Z Z^2_{g_2}\Lambda^d\frac{\Delta^2g_2^2}{2}
        \mbf{k}^2\\
        &+Z^{1/2}\Lambda^{d-1}g_2m^2\delta_{rs}
        -Z^{3/2}\Lambda^{d-1}g_2m^2\delta_{rs}
        \frac{Z_\Delta^2\Delta^2}{8}
        \Lambda^{2-d}\int \frac{d^dk}{(2\pi)^d}\frac{\mbf{k}^2}
        {\left(\mbf{k}^2+m_0^2\right)^2}
    \end{aligned}
\end{equation}
To capture the leading logarithmic contribution we perform dimensiona
regularization and use the following formulae,
\begin{equation}
    \begin{aligned}
        \int\frac{d^dk}{(2\pi)^d}\frac{1}{\mbf{k}^2+m_0^2}
        &=\frac{(m_0)^{d-2}}{(4\pi)^{d/2}}\Gamma\left(1-d/2\right) \\
        \left(\frac{4\pi\Lambda^2}{m_0^2}\right)^{\epsilon/2}
        &=1+\frac{\epsilon}{2}\log\left(\frac{4\pi\Lambda^2}
        {m_0^2}\right)+\mathcal{O}(\epsilon^2) \\
        \Gamma(-n+\epsilon)
        &=\frac{(-1)^n}{n!}
        \left[\frac{1}{\epsilon}+\psi(n+1)\right]
        +\mathcal{O}(\epsilon)
    \end{aligned}
\end{equation}
with $\epsilon=2-d$. Now using the standard renormalization conditions
to make the two point function finite we obtain.
\begin{equation}
    \begin{aligned}
        Z &= 1 + \left(\frac{\Delta^2}{2}\right)\frac{1}{\epsilon}\\
        Z_{g_2} &= 1 - \left(\frac{\Delta^2}{2}\right)\frac{1}{\epsilon}\\
        Z_\Delta &= 1 + \left(\frac{\Delta^2}{4}\right)\frac{1}{\epsilon}\\
    \end{aligned}
\end{equation}
The renormalization of the disorder strength $\Delta$ follow 
straightforwardly by considering,
\begin{equation}
    \begin{aligned}
        Z^2_\Delta\Lambda^\epsilon\Delta^2 &= \Delta_B^2 \\
        \frac{d(Z_\Delta^2\Delta^2)}{d\log\Lambda}
        &=-\epsilon Z^2_\Delta\Delta^2\\
        \frac{d\Delta^2}{d\log\Lambda}
        &=-\frac{\epsilon}{\frac{d\log(Z_\Delta^2\Delta^2)}{d\Delta^2}}\\
        \frac{d\Delta^2}{d\log\Lambda}
        &=-\epsilon\Delta^2 + \frac{\Delta^4}{4}\\
        \frac{d\Delta^2}{d\log\Lambda}
        &=(d-2)\Delta^2 + \frac{\Delta^4}{4}   
    \end{aligned}
\end{equation}
The trivial renormalization of the bulk coupling has already been
mentioned above and the subleading logarithmic correction is unimportant
in this case. An interesting aspect of the calculation is that the 
renormalization was driven by the in plane Goldstone mode fluctuation 
and therefore the results shall be identical for the case of planar 
non-collinear antiferromagnets, i.e the XY case.

\begin{thebibliography}{99}

    \bibitem{imry-ma-75}
        Imry and Ma, 75

    \bibitem{mermin-wagner-66}
        Mermin and Wagner, 66

    \bibitem{aharony-78}
        Aharony, 78

    \bibitem{azaria-92}
        Azaria, Delamotte and Mouhanna, 92

    \bibitem{azaria-93}
        Azaria, Delamotte, Delduc and Jolicoeur, 93

    \bibitem{dombre-88}
        Dombre and Read, 88

    \bibitem{brezin-76}
        Brezin and Zinn-Justin, 76

\end{thebibliography}

\end{document}
